\chapter{Forward Kinematics}
\label{ch:fk}

\section{What FK answers}
Forward kinematics maps known joint coordinates to the tool's position and
orientation. Given $\qvec$, each joint contributes a rigid transform. For a
serial chain, transforms compose by matrix multiplication because the output
frame of one joint becomes the input frame of the next. The result belongs to
$\SE(3)$ and has a rotation matrix and translation vector:
\begin{equation}
 {}^BT_g(\qvec)=
 \begin{bmatrix}{}^BR_g(\qvec)&{}^B\pvec_g(\qvec)\\0&1\end{bmatrix}.
 \label{eq:fk-result}
\end{equation}
The first three columns of the 3-by-3 block are the grasp frame's local unit
axes expressed in base coordinates. The last column is the grasp-frame origin
in metres.

\section{Transform-chain multiplication}
Under the standard-DH construction, FK is \cref{eq:dh-chain}:
\begin{equation}
 {}^BT_g=A_1A_2A_3A_4A_5A_6{}^6T_g.
\end{equation}
One can build intermediate poses
\begin{equation}
 {}^BT_1=A_1,\quad {}^BT_2=A_1A_2,\quad
 {}^BT_3=A_1A_2A_3,\quad {}^BT_6=A_1\cdots A_6,
\end{equation}
then apply the fixed tool transform to obtain ${}^BT_g$. In code, matrix
composition is associative but not commutative: swapping adjacent transforms
changes the reference frame and the physical result.

\section{Fully worked home example}
At $q_i=0$, $v_2=-\pi/2$. The individual matrices reduce to
\begin{align*}
A_1&=\begin{bmatrix}1&0&0&.040\\0&0&1&0\\0&-1&0&.330\\0&0&0&1\end{bmatrix},&
A_2&=\begin{bmatrix}0&-1&0&0\\-1&0&0&-.445\\0&0&-1&0\\0&0&0&1\end{bmatrix},\\
A_3&=\begin{bmatrix}1&0&0&.040\\0&0&1&0\\0&-1&0&0\\0&0&0&1\end{bmatrix},&
A_4&=\begin{bmatrix}1&0&0&0\\0&0&-1&0\\0&1&0&-.440\\0&0&0&1\end{bmatrix},\\
A_5&=\begin{bmatrix}1&0&0&0\\0&0&1&0\\0&-1&0&0\\0&0&0&1\end{bmatrix},&A_6&=I_4.
\end{align*}
Multiplying the first six produces the DH wrist-frame origin $(0.480,0,0.815)$
m and orientation $C_{3\times3}$. Applying ${}^6T_g$ moves the grasp point
approximately $(+0.225980,-0.012300,-0.000911)$ m in base coordinates. Thus
\begin{equation}
 {}^BT_g(0)\approx
 \begin{bmatrix}
 0&0&1&0.705980\\
 0&-1&0&-0.012300\\
 1&0&0&0.814089\\
 0&0&0&1
 \end{bmatrix}.
 \label{eq:home-fk}
\end{equation}
The tiny off-axis entries in direct floating-point URDF evaluation are caused
by the URDF's decimal approximations to $\pi$; \cref{eq:home-fk} is rounded.

\section{Worked nonzero configuration}
The independently recorded verification example
$\qvec=(0.4,0.6,1.2,-0.7,0.8,1.1)$ rad lies within the six arm joint limits.
Applying the URDF chain gives the following rounded pose:
\begin{equation}
{}^BT_g\approx\begin{bmatrix}
-0.857809731&-0.290135590&0.424244982&0.675035771\\
0.156339523&-0.933613380&-0.322372473&0.161569746\\
0.489612519&-0.210207986&0.846222302&1.166928133\\
0&0&0&1
\end{bmatrix}.
\label{eq:fk-nonzero-example}
\end{equation}
The position result is therefore approximately
$(0.675036,0.161570,1.166928)$ m. The rotation's column vectors remain unit
length and mutually orthogonal up to round-off, and its determinant is one.
This is a mathematical FK example, not a measured Gazebo pose.

\section{Executable URDF FK}
\code{gp7_motion/kinematics.py} parses the URDF into a list of transforms and
axes in \code{Kinematics.__init__}. For each chain joint,
\code{fk(q)} first right-multiplies the static joint origin. If the joint is
revolute, it then right-multiplies a rotation generated from the normalized
axis multiplied by its corresponding $q_i$. Fixed joints contribute only
their static transform. Because the chain starts at \code{base_link} and
ends at \code{grasp_frame}, the return value is already the base-to-grasp
transform.

No arm dimensions are duplicated in the runtime FK or IK: changing the URDF
joint origin updates these computations. This reduces parameter drift between
the model and planner, though it means a malformed URDF can also corrupt both
the runtime results and any tests that reuse the same parser.

\section{Independent DH agreement check}
\file{docs/verify_dh.py} independently builds DH matrices using a separately
stated table and reads the exact URDF fixed tool transforms. It checks four
explicit configurations, then 100 deterministic joint vectors drawn uniformly
between the parsed lower and upper limits using NumPy generator seed 42. It
compares translation by Euclidean norm and rotation by the angle of the
relative rotation. In the verification recorded on 14 November 2026, all 104
passed with maximum errors
\begin{equation}
 \max_k\|p_{DH,k}-p_{URDF,k}\|_2=5.50\times10^{-16}\si{m},\qquad
 \max_k\angle(R_{DH,k}^{T}R_{URDF,k})=4.93\times10^{-16}\si{rad}.
\end{equation}
These are machine-precision comparisons of two calculations for the same
implemented URDF. They do not measure physical robot accuracy or CAD fidelity.

\section{Reproduce the FK check}
From the repository root after installing system NumPy/SciPy, use:
\begin{lstlisting}[language=bash]
/usr/bin/python3 docs/verify_dh.py
\end{lstlisting}
The script is read-only. Its current recorded output prints each explicit
configuration's errors and the maximum over all 104. Chapter 13 distinguishes
this mathematical check from the dynamic Gazebo run.
